\begin{textsample}{POS Dim 4 – vem\_ed\_21 – Score 2.00 – vem\_ed\_21/t105.txt}  \label{ex:f4_pos_100}
continuação A equipe dos Projetos Colaborativos Internacionais (PCIs / Cesu) apresentou presencialmente duas \textbf{oficinas}, três comunicações orais e fez quatro moderações, resumidas nas páginas a seguir.
Osvaldo Succi Junior, coordenador dos PCIs / Cesu, compôs o painel de encerramento da conferência, “What is the future of COIL / Intercâmbio Virtual”, com Jon Rubin (COIL Virtual Exchange Foundation / EUA), Mirjam Hauck (The Open University / Inglaterra) e a participação \textbf{remota} de Samia Chasi (International Education Association of South Africa / África do Sul).
A moderação foi de Lavern Samuels (Durban University of Technology / África do Sul).
No painel, Succi Junior enfatizou a importância de adicionar metodologias compatíveis à abordagem do COIL para atender melhor os perfis de professores e as necessidades dos alunos.
Ressaltou que o falante nativo precisa aprender a conversar com seus parceiros internacionais (devagar, sem gírias) e destacou a importância de conectar os PCIs com a comunidade (entorno das Fatecs, prefeituras, empresas).

% matched lemmas: oficina, remoto
\end{textsample}
